\chapter{TINJAUAN PUSTAKA}

\section{Kerangka Teori}
Bagian ini memaparkan teori-teori dasar yang menjadi landasan dalam pengembangan dan analisis sistem transkripsi terdistribusi.

\subsection{Automatic Speech Recognition (ASR) dan Model Whisper}
\foreign{Automatic Speech Recognition} (ASR) adalah teknologi yang memungkinkan komputer untuk mengubah sinyal suara menjadi teks secara otomatis. Perkembangan terbaru dalam bidang ini didominasi oleh pendekatan \foreign{deep learning}, salah satunya adalah model \foreign{Whisper} yang dikembangkan oleh OpenAI. \foreign{Whisper} menggunakan arsitektur \foreign{Transformer} berbasis \foreign{encoder-decoder} yang dilatih menggunakan 680.000 jam data audio multibahasa secara \foreign{weakly supervised} \cite{radford_robust_nodate}. Hal ini membuat model tersebut memiliki tingkat ketahanan (\foreign{robustness}) yang tinggi terhadap aksen, kebisingan latar belakang, dan istilah teknis dibandingkan model ASR tradisional.

\subsection{Sistem Transkripsi Terdistribusi}
Sistem terdistribusi adalah sekumpulan komputer mandiri yang tampak bagi pengguna sebagai satu sistem tunggal yang koheren. Dalam konteks transkripsi otomatis, beban kerja pemrosesan audio bersifat \foreign{compute-intensive}. Oleh karena itu, penggunaan arsitektur terdistribusi yang memisahkan antara antarmuka pengguna (\foreign{Next.js}) dan unit pemrosesan (\foreign{Python Worker}) sangat krusial. Komunikasi antar komponen ini biasanya dilakukan melalui \foreign{Message Broker} seperti Redis atau RabbitMQ untuk mengelola antrean tugas secara asinkron.

\subsection{Metrik Evaluasi: Word Error Rate (WER)}
Akurasi sistem ASR diukur menggunakan metrik \foreign{Word Error Rate} (WER). WER dihitung dengan menjumlahkan jumlah substitusi ($S$), penghapusan ($D$), dan penyisipan ($I$) kata, kemudian dibagi dengan jumlah total kata ($N$) pada referensi asli:
\begin{equation}
    WER = \frac{S + D + I}{N}
\end{equation}
Nilai WER yang lebih rendah menunjukkan tingkat akurasi yang lebih tinggi.

\section{Literatur Review}
Penelitian mengenai implementasi dan optimasi model \foreign{Whisper} telah banyak dilakukan dalam beberapa tahun terakhir. Fathony dkk. \cite{fathony_optimisasi_nodate} melakukan penelitian mengenai optimasi \foreign{Whisper} untuk Bahasa Indonesia menggunakan pendekatan \foreign{Hybrid Cloud}. Temuan mereka menunjukkan bahwa penggunaan strategi \foreign{multi-engine} dapat meningkatkan kecepatan transkripsi secara signifikan tanpa mengorbankan akurasi secara drastis pada data bahasa lokal.

Penelitian lain oleh Andreyev \cite{andreyev_quantization_2025} berfokus pada teknik kuantisasi model \foreign{Whisper} (INT4, INT5, dan INT8) untuk mengurangi latensi dan penggunaan memori pada perangkat \foreign{edge}. Hasilnya menunjukkan bahwa kuantisasi mampu mengurangi penggunaan memori hingga 45\% dengan penurunan akurasi yang minimal. Hal ini relevan dengan penelitian ini yang juga mempertimbangkan penggunaan sumber daya GPU yang terbatas di lingkungan kampus.

Bain dkk. \cite{bain_whisperx_2023} memperkenalkan \foreign{WhisperX} yang mengoptimalkan model \foreign{Whisper} asli dengan menambahkan fitur \foreign{Voice Activity Detection} (VAD) dan \foreign{forced alignment} untuk menghasilkan \foreign{timestamp} tingkat kata yang lebih akurat. Meskipun penelitian ini berfokus pada model \foreign{Whisper} standar, teknik pemrosesan yang diusulkan dalam \foreign{WhisperX} menjadi referensi penting dalam memahami efisiensi pemrosesan audio berdurasi panjang.

Perbedaan utama penelitian ini dengan penelitian sebelumnya adalah pada fokus implementasi sistem transkripsi yang terintegrasi secara penuh dalam ekosistem SaaS terdistribusi menggunakan \foreign{Next.js} dan \foreign{Python Worker}, serta melakukan \foreign{benchmarking} komprehensif pada seluruh varian model \foreign{Whisper} menggunakan data rekaman lokal bahasa Indonesia guna mencari titik temu antara efisiensi biaya infrastruktur dan akurasi hasil.

\section{Hipotesis}
Berdasarkan tinjauan pustaka di atas, diajukan hipotesis sebagai berikut:
\begin{enumerate}
    \item Terdapat korelasi positif antara ukuran model \foreign{Whisper} dengan tingkat akurasi, di mana model \foreign{Large-v3} diharapkan menghasilkan WER terendah namun dengan \foreign{inference time} paling lama.
    \item Model \foreign{Large-v3-Turbo} mampu memberikan akurasi yang mendekati model \foreign{Large-v3} dengan kecepatan pemrosesan yang jauh lebih efisien pada infrastruktur NVIDIA RTX A4000.
    \item Arsitektur terdistribusi menggunakan \foreign{Python Worker} mampu menjaga stabilitas sistem transkripsi pada beban kerja tinggi dibandingkan dengan pemrosesan monolitik.
\end{enumerate}
